\section*{v2.7.0 Field-Energy Realization and Semigroup Readout Core}
\addcontentsline{toc}{section}{v2.7.0 Field-Energy Realization and Semigroup Readout Core}

\begin{quote}
\textbf{Normalization note (v2.17.1 local review).}
This block is retained as a \emph{historical realization and semigroup-readout precursor}. It supports the later field/probe realization layer and should not be read as a second independent primary core.
\end{quote}

The next structural step after forward-invariant ledger closure is to show that the
weighted ledger is not only a bookkeeping device but already the shadow of a genuine
field-energy balance. Version~\docversion\ therefore adds a realization layer: the
channel budgets are interpreted as coarse observables of an underlying evolution
system, and admissibility is reformulated as dissipative semigroup compatibility.
This still does not claim the final field equations, but it fixes the exact kind of
continuous dynamics those equations must generate if they are to remain faithful to
the structural proof architecture.

\subsection*{1. Abstract state space and energy readout}
Let
\[
X := X_{\mathrm{tri}}\oplus X_{\perp}\oplus X_{\mathrm{meas}}\oplus X_{\mathrm{res}}
\]
be an abstract Banach or Hilbert state space carrying the channel-resolved structural
variables. We write a state as
\[
U=(u_{\mathrm{tri}},u_{\perp},u_{\mathrm{meas}},u_{\mathrm{res}}).
\]
The ledger variables are no longer treated as primitive numbers but as readout
functionals
\[
B_i(U)\ge 0,
\qquad i\in\{\mathrm{tri},\perp,\mathrm{meas},\mathrm{res}\},
\]
with the weighted energy
\[
\mathcal E(U):=\alpha_{\mathrm{tri}}B_{\mathrm{tri}}(U)
+\alpha_{\perp}B_{\perp}(U)
+\alpha_{\mathrm{meas}}B_{\mathrm{meas}}(U)
+\alpha_{\mathrm{res}}B_{\mathrm{res}}(U).
\]
The content of this step is that $\mathcal E$ should reproduce the Lyapunov ledger
$\mathcal L$ up to uniform equivalence on the admissible sector. Thus the ledger is
interpreted as a compressed energy observable of the underlying state rather than an
external annotation.

\subsection*{2. Dissipative generator Ansatz}
Assume that the channel evolution is generated by an abstract equation
\[
\partial_\tau U = \mathcal A U + \mathcal N(U) + \mathcal S_{\mathrm{tag}}(\tau),
\]
where $\mathcal A$ is the linear transport-dissipation generator, $\mathcal N$ is a
structured nonlinear coupling compatible with HC confinement, and
$\mathcal S_{\mathrm{tag}}$ collects the tagged residual sources. The admissibility
requirement from the earlier sections now becomes the statement that
$\mathcal A+\mathcal N$ is energy-dissipative modulo tagged production and match
error, namely
\[
\frac{d}{d\tau}\mathcal E(U(\tau))
\le -c_0\mathcal E(U(\tau))
+ C_{\mathrm{tag}}\,\mathrm{Res}_{\mathrm{tag}}(\tau)
+ C_{\mathrm{match}}\,\Delta_{\mathrm{match}}(\tau)
\]
for every admissible trajectory. This is the semigroup form of the weighted ledger
closure derived in v2.6.0.

\subsection*{3. Realization principle}
The structural point can now be isolated as a single principle.

\paragraph{Principle 1 (field-energy realization).}
An eventual field equation is acceptable only if its induced state evolution admits a
channel readout $U\mapsto (B_i(U))_i$ and a weighted energy $\mathcal E$ such that:
\begin{enumerate}[label=(\alph*)]
  \item $\mathcal E$ is uniformly equivalent to the admissible ledger functional $\mathcal L$;
  \item the induced evolution preserves the forward-invariant admissible sector;
  \item all positive production terms are either dissipatively absorbed or explicitly tagged; and
  \item probe mismatch enters only through a controlled source term $\Delta_{\mathrm{match}}$.
\end{enumerate}
Any ansatz failing one of these four points is structurally too weak, even if it
looks formally plausible at equation level.

\subsection*{4. Semigroup consequence and proof memory}
If the above dissipative estimate holds, then the evolution operator $S(\tau)$
generated by the admissible dynamics satisfies an estimate of the form
\[
\mathcal E(S(\tau)U_0)
\le e^{-c_0\tau}\mathcal E(U_0)
+\int_0^\tau e^{-c_0(\tau-s)}
\bigl(C_{\mathrm{tag}}\,\mathrm{Res}_{\mathrm{tag}}(s)
+ C_{\mathrm{match}}\,\Delta_{\mathrm{match}}(s)\bigr)\,ds.
\]
Hence proof memory can now be reformulated in semigroup language: restarting from a
checkpoint is valid because the future evolution depends continuously on the current
energy state and the tagged source history, not on the unrecoverable microscopic path.
This upgrades the earlier restart argument into an abstract continuous-time stability
statement.

\subsection*{5. Probe-sector realization}
The same realization principle applies to the two-probe geometry. If
\[
U_{\mathrm{probe}}=(U^{(1)},U^{(2)},U^{(Q)},U^{(M)},U^{(R)})
\]
denotes the probe state and $\mathcal E_{\mathrm{probe}}$ the corresponding weighted
energy, then admissibility requires
\[
\frac{d}{d\tau}\mathcal E_{\mathrm{probe}}
\le -c_{\mathrm{probe}}\mathcal E_{\mathrm{probe}}
+ C^{\mathrm{probe}}_{\mathrm{tag}}\,\mathrm{Res}^{\mathrm{probe}}_{\mathrm{tag}}
+ C^{\mathrm{probe}}_{\mathrm{match}}\,\Delta_{\mathrm{match}}.
\]
Thus the probe mechanism is now seen as a dissipative comparison flow living in the
same semigroup class as the main derivation. This makes the search geometry more than
a metaphor: it becomes an admissible evolution subsystem with inherited energy
control.

\subsection*{6. Consequence for future field equations and appendix tools}
The field-equation programme is sharpened once again. A candidate PDE or abstract
operator equation must now be checked in the following order:
\begin{enumerate}[label=(\arabic*)]
  \item identify the state space and channel decomposition;
  \item construct the ledger readout functionals $B_i(U)$;
  \item verify equivalence between readout energy $\mathcal E$ and structural ledger $\mathcal L$;
  \item prove dissipativity modulo tagged sources; and
  \item verify probe inheritance and restart compatibility.
\end{enumerate}
Accordingly, appendix tools can receive one further classification layer. A strong
tool may now be marked as a state-space constructor, energy comparator, semigroup
stability lemma, or source-tagging certificate. This helps separate tools that only
support symbolic rearrangement from those that certify genuine dynamic admissibility.

\subsection*{7. v\docversion\ readiness criterion}
The present step should be regarded as successful if the following can now be done:
\begin{itemize}
  \item the ledger is interpretable as a uniformly equivalent field-energy observable;
  \item admissibility is reformulated as dissipative semigroup compatibility;
  \item proof memory is supported by a continuous-time stability estimate;
  \item the two-probe sector is embedded into the same realization framework; and
  \item future field equations can be accepted or rejected by an explicit realization checklist.
\end{itemize}

This is the gain of v\docversion. The companion monolith now carries not only
closure, channel specialization, transport dynamics, and forward invariance, but also
a first abstract field-energy realization principle. The later field-equation layer is
therefore no longer free-form: it has to generate the same dissipative proof memory,
probe inheritance, and tagged-source discipline already fixed at the structural level.


